% Source: https://artofproblemsolving.com/community/c3h2125065_400_posts_lets_go
% Stable identity: aops-topic-2125065-q01
A student ``proves'' that $\dfrac{x+y}{2}\ge\sqrt{xy}$ as follows. By the trivial inequality, $(x-y)^2\ge0$, so $x^2-2xy+y^2\ge0$. Adding $4xy$ to both sides gives $x^2+2xy+y^2\ge4xy$, and dividing by $4$ gives $\dfrac{(x+y)^2}{4}\ge xy$. Taking square roots of both sides gives $\dfrac{x+y}{2}\ge\sqrt{xy}$.
\begin{enumerate}
\item[(a)] Show that the conclusion is false for some real numbers $x$ and $y$ for which $\sqrt{xy}$ is defined.
\item[(b)] Which step is not justified in general? Correct the argument to prove that $\dfrac{x+y}{2}\ge\sqrt{xy}$ for all nonnegative real numbers $x$ and $y$, and find when equality holds.
\end{enumerate}
